% !TeX root = ../../Book3.tex
% 纲要标题：### C.1 卓越性、形式纤维与普适链性质

\section{卓越性、形式纤维与普适链性质}
\label{b3:sec:C01}

% 纲要细目（与计划纲要同步）：
% * 卓越性、形式纤维与普适链性质的动机
%
% <!-- 短节：不设小节 -->


% 正文以具体模型说明专题；长证明给出概要并指明技术细节的出处。

对一个局部环取完备化，常能把多项式计算改成幂级数计算。但忠实平坦本身并没有说明每条纤维的几何性质，也没有保证所有有限型代数的正则点都组成开集。研究奇点时，往往还希望完成下列工作：在完备化前后比较局部结构；在一个族中找出统一的正则开集；沿不同素理想链计算一致的余维。卓越性把这些需求组织为明确条件。

\ProseParagraphBreak
设 \((A,\mathfrak m)\) 为 Noether 局部环，\(\widehat A\) 为其 \(\mathfrak m\) 进完备化。对 \(\mathfrak p\in\operatorname{Spec}A\)，环
\[
\widehat A\otimes_A\kappa(\mathfrak p)
\]
称为该素点处的\term{形式纤维环}[formal fibre ring]。它不是只考察闭点 \(\mathfrak m\)：闭纤维确实为 \(A/\mathfrak m\)，其他素点的纤维则可能包含新的信息。相关定义见松村英之《Commutative Ring Theory》\cite[\S 32, pp. 255--256]{Matsumura1989CommutativeRingTheory}。

\begin{example}\label{b3:exm:formal-fibres-line}
设 \(k\) 为特征零域，\(A=k[x]_{(x)}\)，则 \(\widehat A=k[[x]]\)。\(A\) 的两个素点 \((x)\) 与 \((0)\) 上的形式纤维分别为
\[
k[[x]]\otimes_A k\cong k,
\qquad
k[[x]]\otimes_A k(x)\cong k((x)).
\]
第二式中，对 \(A\) 所有非零元局部化，等价于使 \(x\) 可逆：非零多项式在 \(k[[x]]\) 中是某个 \(x\) 次幂乘单位，反过来 \(x\) 本来就在局部化集合中。于是得到 Laurent 级数域，而不是把泛纤维误认成剩余域 \(k\)。这两个纤维都是几何正则的：在特征零域上作任意有限扩域都是可分扩张，与一个域张量后为有限个域之积，故仍正则。
\end{example}

“几何正则”要求在有限基域扩张后仍为正则环。要求 \(A_{\mathfrak p}\to\widehat{A_{\mathfrak p}}\) 的所有纤维都几何正则，便得到 G 环条件。另一个条件是普适链性质：每个有限型 \(A\) 代数中，从一个固定素理想到另一个固定素理想的饱和链具有相同长度。第三个条件通常称为 J-2，要求每个有限型 \(A\) 代数的正则点集为开集。\cref{b3:def:excellent-ring}统一给出这三个条件及卓越环的定义；此处强调它们控制的是不同问题。

\ProseParagraphBreak
域上的有限型代数、它们的局部化以及完备 Noether 局部环属于卓越环；\cref{b3:thm:excellent-basic-classes}给出陈述及证明概要，分别说明链性质、形式纤维及正则点开放性的验证。例如尖点的局部环虽然有奇点，仍是卓越环。因此卓越性并不表示“没有奇点”；它给出研究奇点时所需的一组良好比较性质。读者在具体问题中调用它时，仍应指明究竟需要形式纤维、正则点的开放性，还是链长度条件。
